\chapter{Modelo ambiental}
\label{cap:ambiental}

El modelo ambiental define la frontera entre el sistema y su entorno: qué hace el sistema,
con quién se comunica y a qué estímulos debe responder \parencite{yourdon1989}. Se compone de
la declaración de propósitos, el diagrama de contexto y la lista de acontecimientos.

% =========================================================
\section{Declaración de propósitos}
\label{sec:propositos}

El propósito de BeeTracer es \textbf{registrar de forma verificable el estado de la carga en
cada etapa crítica del desconsolidado de contenedores}, desde la llegada del manifiesto
electrónico hasta la entrega del reporte al cliente, para que el terminal extraportuario pueda
demostrar, ante un reclamo o un seguro, si la mercancía venía dañada desde origen o si el
daño ocurrió durante sus propias maniobras.

Para cumplir ese propósito, el sistema:

\begin{enumerate}
    \item Conoce el contenido de cada contenedor antes de su arribo, importando el manifiesto
    electrónico del Servicio Nacional de Aduanas, y admite el ingreso manual de la carga
    (tarja ciega) cuando el manifiesto no está disponible.
    \item Planifica cada faena de desconsolidado, asignando fecha, franja horaria, cuadrilla y
    jefe tarjador sin conflictos de agenda.
    \item Exige evidencia fotográfica obligatoria, fechada y geolocalizada, del sello, la
    apertura, cada lote separado por cliente, cada incidencia y el contenedor vacío.
    \item Contrasta lo recibido con lo declarado y registra las incidencias (daño, faltante,
    sobrante, mal embalaje y sello alterado) asociadas al cliente afectado.
    \item Somete el reporte de cada operación a la validación de un supervisor antes de que
    llegue al cliente, permitiendo rechazar y recapturar la evidencia defectuosa.
    \item Entrega al cliente el reporte aprobado y le permite consultar el estado de sus
    operaciones, con acceso restringido a la carga consignada a su empresa.
\end{enumerate}

Quedan fuera del propósito del sistema la manipulación física de la carga, el cobro de los
servicios del depósito, la gestión del patio de contenedores (visor 3D) y los avisos de
transporte, que Bilix comercializa como módulos separados.

% =========================================================
\section{Diagrama de contexto}
\label{sec:contexto}

El diagrama de contexto (Figura~\ref{fig:contexto}) representa a BeeTracer como un único
proceso, numerado 0, rodeado por los \emph{terminadores}: las personas y sistemas externos que
envían o reciben datos. Se identifican seis terminadores, que corresponden a los actores del
Informe 2:

\begin{itemize}
    \item \textbf{Administrador del depósito}: envía las solicitudes de mantención de los datos
    maestros y recibe su resultado.
    \item \textbf{Jefe de faena}: solicita la importación de manifiestos, programa las faenas,
    consulta las operaciones y solicita el reenvío de reportes.
    \item \textbf{Servicio Nacional de Aduanas}: sistema externo que recibe la solicitud de
    manifiesto, entrega el manifiesto XML y recibe la notificación del resultado de la
    importación.
    \item \textbf{Jefe tarjador}: se identifica en la aplicación móvil, recibe sus faenas
    asignadas, envía el registro de terreno (datos y fotografías) y, si se le solicita,
    recaptura la evidencia rechazada.
    \item \textbf{Supervisor del reporte}: solicita la revisión de los reportes, recibe su
    contenido, decide su aprobación o rechazo y recibe el resultado del envío.
    \item \textbf{Cliente del depósito}: consulta el estado de sus operaciones y recibe el
    reporte PDF aprobado.
\end{itemize}

La cuadrilla no aparece como terminador porque no intercambia datos con el sistema: su
trabajo se registra a través del jefe tarjador. Tampoco aparece el servidor de correo, ya que
es el \emph{medio} por el cual el reporte llega al cliente y no un origen o destino de datos;
en el modelo esencial el flujo \este{reporte PDF} va directamente al Cliente del depósito.

\diagramatikz{contexto}
    {Diagrama de contexto}
    {Diagrama de contexto del sistema BeeTracer.}
    {fig:contexto}

% =========================================================
\section{Lista de acontecimientos}
\label{sec:acontecimientos}

La lista de acontecimientos enumera los sucesos del entorno ante los cuales el sistema debe
responder. Cada acontecimiento se clasifica según \textcite{yourdon1989} como de
\textbf{flujo (F)}, cuando llega un flujo de datos al sistema; \textbf{temporal (T)}, cuando
lo dispara el paso del tiempo; o de \textbf{control (C)}, cuando lo dispara una señal sin
datos. Todos los acontecimientos de BeeTracer son de flujo: el sistema siempre reacciona a
datos que envía un terminador. El único candidato a acontecimiento de control, la
recuperación de la señal en el patio, pertenece a la tecnología de implementación y por ello
no forma parte del modelo esencial (se trata en el diseño, Capítulo~\ref{cap:estructura}).

La Tabla~\ref{tab:acontecimientos} indica, para cada acontecimiento, el flujo de entrada que lo
anuncia y el proceso del diagrama de flujo de datos de nivel~1 que le responde.

{\estilotabla
\begin{longtable}{|C{0.6cm}|L{6.2cm}|C{0.9cm}|L{3.4cm}|C{1.9cm}|}
\caption{Lista de acontecimientos del sistema BeeTracer.}
\label{tab:acontecimientos}\\
\hline
\rowcolor{filacabecera}
\textbf{N.º} & \textbf{Acontecimiento} & \textbf{Tipo} & \textbf{Flujo de entrada} &
\textbf{Proceso} \\
\hline
\endfirsthead
\hline
\rowcolor{filacabecera}
\textbf{N.º} & \textbf{Acontecimiento} & \textbf{Tipo} & \textbf{Flujo de entrada} &
\textbf{Proceso} \\
\hline
\endhead
1 & El Administrador solicita crear, editar, eliminar o cambiar el estado de un usuario. &
F & solicitud de mantención (datos de usuario) & 1.1 \\ \hline
2 & El Administrador solicita mantener un rol y los permisos que otorga. &
F & solicitud de mantención (datos de rol) & 1.2 \\ \hline
3 & El Administrador solicita mantener un tipo de bulto. &
F & solicitud de mantención (datos de tipo de bulto) & 1.3 \\ \hline
4 & El Administrador solicita mantener un cliente y su contacto. &
F & solicitud de mantención (datos de cliente) & 1.4 \\ \hline
5 & El Administrador solicita mantener un trabajador, una cuadrilla o la asignación de un
trabajador a una cuadrilla. &
F & solicitud de mantención (datos de trabajador) & 1.5 \\ \hline
6 & El Jefe de faena solicita importar el manifiesto de una nave y viaje. &
F & solicitud de importación & 2.1 \\ \hline
7 & El Servicio Nacional de Aduanas entrega el manifiesto solicitado. &
F & manifiesto XML & 2.2 \\ \hline
8 & El Jefe de faena crea la programación del desconsolidado de un contenedor. &
F & datos de programación (nueva programación) & 3.1 \\ \hline
9 & El Jefe de faena asigna la cuadrilla y el jefe tarjador de una programación. &
F & datos de programación (asignación) & 3.2 \\ \hline
10 & El Jefe de faena reprograma, anula, elimina o cambia el estado de una programación. &
F & datos de programación (cambio de programación) & 3.3 \\ \hline
11 & El Jefe tarjador inicia sesión en la aplicación móvil y solicita sus faenas asignadas. &
F & credenciales & 3.4 \\ \hline
12 & El Jefe tarjador inicia la operación de una faena asignada. &
F & registro de terreno (datos de inicio) & 4.1 \\ \hline
13 & El Jefe tarjador registra el sello del contenedor y su fotografía. &
F & registro de terreno (datos de sello) & 4.2 \\ \hline
14 & El Jefe tarjador registra un lote separado por cliente y su fotografía. &
F & registro de terreno (datos de lote) & 4.3 \\ \hline
15 & El Jefe tarjador registra una incidencia con sus fotografías de respaldo. &
F & registro de terreno (datos de incidencia) & 4.4 \\ \hline
16 & El Jefe tarjador cierra la faena con la fotografía del contenedor vacío. &
F & registro de terreno (datos de cierre) & 4.5 \\ \hline
17 & El Supervisor del reporte solicita revisar un reporte pendiente. &
F & solicitud de revisión & 5.2 \\ \hline
18 & El Supervisor del reporte aprueba un reporte. &
F & decisión de validación (aprobación) & 5.3 \\ \hline
19 & El Supervisor del reporte rechaza una fotografía defectuosa y registra una observación. &
F & decisión de validación (rechazo) & 5.4 \\ \hline
20 & El Jefe tarjador envía la fotografía que reemplaza a la rechazada. &
F & fotografía de reemplazo & 5.5 \\ \hline
21 & El Jefe de faena consulta las operaciones o modifica el estado de una de ellas. &
F & consulta de operaciones & 6.2 \\ \hline
22 & El Jefe de faena solicita reenviar al cliente un reporte aprobado. &
F & solicitud de reenvío & 6.1 \\ \hline
23 & El Cliente del depósito consulta el estado de sus operaciones. &
F & consulta de estado & 6.3 \\ \hline
24 & El Cliente del depósito solicita visualizar o descargar un reporte aprobado. &
F & solicitud de PDF & 6.4 \\ \hline
\end{longtable}
}

Algunas respuestas del sistema no tienen un acontecimiento propio porque son consecuencia de
otro: al cerrarse la faena (acontecimiento 16) se genera el borrador del reporte (proceso 5.1)
y al aprobarse un reporte (acontecimiento 18) se envía al cliente (proceso 6.1). Estas cadenas
se representan en los diagramas de flujo de datos mediante los flujos internos
\este{operación finalizada} y \este{reporte aprobado}.
